%%%%%%%%%%%%%%%%
%%% deut 3 %%%
%%%%%%%%%%%%%%%%

\Chapter{3}

\Verse{1}
{
	\hDtr1{וַנֵּפֶן וַנַּעַל דֶּרֶךְ הַבָּשָׁן וַיֵּצֵא עוֹג מֶלֶךְ הַבָּשָׁן לִקְרָאתֵנוּ הוּא וְכל עַמּוֹ לַמִּלְחָמָה אֶדְרֶעִי׃}
}
{
	\eDtr1{
		3 “And we turned and went up the way to Bashan. And Og,\\
		king of Bashan, came out at us, he and all his people, for\\
		war at Edrei.\\
	}
}
{
}

\Verse{2}
{
	\hDtr1{וַיֹּאמֶר יְהֹוָה אֵלַי אַל תִּירָא אֹתוֹ כִּי בְיָדְךָ נָתַתִּי אֹתוֹ וְאֶת כּל עַמּוֹ וְאֶת אַרְצוֹ וְעָשִׂיתָ לּוֹ כַּאֲשֶׁר עָשִׂיתָ לְסִיחֹן מֶלֶךְ הָאֱמֹרִי אֲשֶׁר יוֹשֵׁב בְּחֶשְׁבּוֹן׃}
}
{
	\eDtr1{
		And YHWH said to me, ‘Don’t be afraid of him, because I’ve\\
		put him and all his people and his land in your hand, and\\
		you shall do to him as you did to Sihon, king of the\\
		Amorites, who lived in Heshbon.’\\
	}
}
{
}

\Verse{3}
{
	\hDtr1{וַיִּתֵּן יְהֹוָה אֱלֹהֵינוּ בְּיָדֵנוּ גַּם אֶת עוֹג מֶלֶךְ הַבָּשָׁן וְאֶת כּל עַמּוֹ וַנַּכֵּהוּ עַד בִּלְתִּי הִשְׁאִיר לוֹ שָׂרִיד׃}
}
{
	\eDtr1{
		And YHWH, our God, put Og, king of Bashan, and all his\\
		people in our hand as well, and we struck him until he did\\
		not have a remnant left.\\
	}
}
{
}

\Verse{4}
{
	\hDtr1{וַנִּלְכֹּד אֶת כּל עָרָיו בָּעֵת הַהִוא לֹא הָיְתָה קִרְיָה אֲשֶׁר לֹא לָקַחְנוּ מֵאִתָּם שִׁשִּׁים עִיר כּל חֶבֶל אַרְגֹּב מַמְלֶכֶת עוֹג בַּבָּשָׁן׃}
}
{
	\eDtr1{
		And we captured all his cities at that time. There wasn’t\\
		a town that we didn’t take from them: sixty cities—all the\\
		region of Argob, Og’s kingdom in Bashan.\\
	}
}
{
}

\Verse{5}
{
	\hDtr1{כּל אֵלֶּה עָרִים בְּצֻרֹת חוֹמָה גְבֹהָה דְּלָתַיִם וּבְרִיחַ לְבַד מֵעָרֵי הַפְּרָזִי הַרְבֵּה מְאֹד׃}
}
{
	\eDtr1{
		All of these were ortified cities—a high wall, double\\
		gates, and a bar—aside from a great many unwalled cities.\\
	}
}
{
}

\Verse{6}
{
	\hDtr1{וַנַּחֲרֵם אוֹתָם כַּאֲשֶׁר עָשִׂינוּ לְסִיחֹן מֶלֶךְ חֶשְׁבּוֹן הַחֲרֵם כּל עִיר מְתִם הַנָּשִׁים וְהַטָּף׃}
}
{
	\eDtr1{
		And we completely destroyed them as we did to Sihon, king\\
		of Heshbon, completely destroying every city: men, women,\\
		and infants.\\
	}
}
{
%	\footnote{
%		completely destroyed. See the comment on Deut 2:34.
%	}
}

\Verse{7}
{
	\hDtr1{וְכל הַבְּהֵמָה וּשְׁלַל הֶעָרִים בַּזּוֹנוּ לָנוּ׃}
}
{
	\eDtr1{
		And we despoiled all the animals and the spoil of the\\
		cities for ourselves.\\
	}
}
{
}

\Verse{8}
{
	\hDtr1{וַנִּקַּח בָּעֵת הַהִוא אֶת הָאָרֶץ מִיַּד שְׁנֵי מַלְכֵי הָאֱמֹרִי אֲשֶׁר בְּעֵבֶר הַיַּרְדֵּן מִנַּחַל אַרְנֹן עַד הַר חֶרְמוֹן׃}
}
{
	\eDtr1{
		And at that time we took the land from the hand of the two\\
		kings of the Amorites that was across the Jordan, from the\\
		Wadi Arnon to Mount Hermon\\
	}
}
{
%	\footnote{
%		across the Jordan. This phrase occurs several times in Moses’ speech in
%		Deuteronomy. Ibn Ezra hinted at a secret implied by this and several
%		other matters in the Torah, and he added, “One who understands should
%		keep silent.” But scholars of later centuries no longer kept silent. The
%		issue, presumably, was that the land in question is across the Jordan
%		only from the point of view of someone writing in Israel. Moses, who
%		never set foot in Israel, would not be expected to refer to the place
%		where he was standing as across the Jordan. Of course we might say that
%		Moses says this phrase with a future audience in mind, a people settled
%		in Israel. But Ibn Ezra’s comments are important as being among the
%		first hints that some of the traditional rabbinic commentators
%		questioned whether Moses wrote the Torah. That question—how the Torah
%		came to exist—has become a central concern of biblical scholarship in
%		the last two centuries. I have written about it elsewhere, but it is not
%		the subject of this commentary. Those who want to pursue it can look at
%		my Who Wrote the Bible? and The Hidden Book in the Bible. Those books
%		look at the persons and events that gave birth to the Torah. This
%		commentary is focused on the Torah itself, the work that those persons
%		and events produced.
%	}
}

\Verse{9}
{
	\hDtr1{צִידֹנִים יִקְרְאוּ לְחֶרְמוֹן שִׂרְיֹן וְהָאֱמֹרִי יִקְרְאוּ לוֹ שְׂנִיר׃}
}
{
	\eDtr1{
		(Sidonians call Hermon Sirion, and the Amorites call it\\
		Senir),\\
	}
}
{
}

\Verse{10}
{
	\hDtr1{כֹּל עָרֵי הַמִּישֹׁר וְכל הַגִּלְעָד וְכל הַבָּשָׁן עַד סַלְכָה וְאֶדְרֶעִי עָרֵי מַמְלֶכֶת עוֹג בַּבָּשָׁן׃}
}
{
	\eDtr1{
		all the cities of the plain and all of Gilead and all of\\
		Bashan as far as Salcah and Edrei, cities of Og’s kingdom\\
		in Bashan.\\
	}
}
{
}

\Verse{11}
{
	\hDtr1{כִּי רַק עוֹג מֶלֶךְ הַבָּשָׁן נִשְׁאַר מִיֶּתֶר הָרְפָאִים הִנֵּה עַרְשׂוֹ עֶרֶשׂ בַּרְזֶל הֲלֹה הִוא בְּרַבַּת בְּנֵי עַמּוֹן תֵּשַׁע אַמּוֹת ארְכָּהּ וְאַרְבַּע אַמּוֹת רחְבָּהּ בְּאַמַּת אִישׁ׃}
}
{
	\eDtr1{
		Because only Og, king of Bashan, was left from the rest of\\
		the Rephaim. Here, his bedstead was a bedstead of iron.\\
		Isn’t it in Rabbah of the children of Ammon? Its length is\\
		nine cubits and its width is four cubits by a man’s cubit.\\
	}
}
{
%	\footnote{
%		his bedstead was a bedstead of iron. Og is pictured as a giant,
%		requiring a tremendous bed that had to made of iron in order to bear his
%		weight. This may also relate to his being the last of the Rephaim, being
%		connected with beings who are now associated with the world of the dead.
%		(See the comment on 2:11.) Footnote 3:11. by a man’s cubit. A cubit is
%		measured from a man’s elbow to his hand’s second knuckle: eighteen
%		inches.
%	}
}

\Verse{12}
{
	\hDtr1{וְאֶת הָאָרֶץ הַזֹּאת יָרַשְׁנוּ בָּעֵת הַהִוא מֵעֲרֹעֵר אֲשֶׁר עַל נַחַל אַרְנֹן וַחֲצִי הַר הַגִּלְעָד וְעָרָיו נָתַתִּי לָראוּבֵנִי וְלַגָּדִי׃}
}
{
	\eDtr1{
		And we took possession of this land at that time. I gave\\
		the Reubenites and the Gadites from Aroer, which is on the\\
		Wadi Arnon, and half of the hill country of Gilead and its\\
		cities.\\
	}
}
{
}

\Verse{13}
{
	\hDtr1{וְיֶתֶר הַגִּלְעָד וְכל הַבָּשָׁן מַמְלֶכֶת עוֹג נָתַתִּי לַחֲצִי שֵׁבֶט הַמְנַשֶּׁה כֹּל חֶבֶל הָאַרְגֹּב לְכל הַבָּשָׁן הַהוּא יִקָּרֵא אֶרֶץ רְפָאִים׃}
}
{
	\eDtr1{
		And I gave the rest of Gilead and all of Bashan, Og’s\\
		kingdom, to half of the tribe of Manasseh. All the region\\
		of Argob: all of that Bashan is called ‘the land of the\\
		Rephaim.’\\
	}
}
{
}

\Verse{14}
{
	\hDtr1{יָאִיר בֶּן מְנַשֶּׁה לָקַח אֶת כּל חֶבֶל אַרְגֹּב עַד גְּבוּל הַגְּשׁוּרִי וְהַמַּעֲכָתִי וַיִּקְרָא אֹתָם עַל שְׁמוֹ אֶת הַבָּשָׁן חַוֺּת יָאִיר עַד הַיּוֹם הַזֶּה׃}
}
{
	\eDtr1{
		Jair son of Manasseh had taken all the region of Argob as\\
		far as the border of the Geshurites nd the Maachathites,\\
		and he called them, the Bashan, by his name, Havvoth-Jair,\\
		to this day.\\
	}
}
{
%	\footnote{
%		Geshurites. Later in the Tanak, Geshur becomes extremely important, a
%		Canaanite people who reside close to the Israelites. King David marries
%		Maacah, a Geshurite princess, and they have a son, Absalom. Absalom
%		kills his elder half-brother, Amnon (for raping Absalom’s full sister,
%		Tamar), who would have been first in line to succeed their father,
%		David. And later Absalom leads a rebellion to usurp his father’s throne
%		(2 Sam 3:3; 13:1–19:1).
%	}
}

\Verse{15}
{
	\hDtr1{וּלְמָכִיר נָתַתִּי אֶת הַגִּלְעָד׃}
}
{
	\eDtr1{And I gave Gilead to Machir.}
}
{
}

\Verse{16}
{
	\hDtr1{וְלָראוּבֵנִי וְלַגָּדִי נָתַתִּי מִן הַגִּלְעָד וְעַד נַחַל אַרְנֹן תּוֹךְ הַנַּחַל וּגְבֻל וְעַד יַבֹּק הַנַּחַל גְּבוּל בְּנֵי עַמּוֹן׃}
}
{
	\eDtr1{
		And I gave the Reubenites and the Gadites from Gilead to\\
		the Wadi Arnon, the middle of the wadi being the border,\\
		and to the Wadi Jabbok, the border of the children of\\
		Ammon,\\
	}
}
{
}

\Verse{17}
{
	\hDtr1{וְהָעֲרָבָה וְהַיַּרְדֵּן וּגְבֻל מִכִּנֶּרֶת וְעַד יָם הָעֲרָבָה יָם הַמֶּלַח תַּחַת אַשְׁדֹּת הַפִּסְגָּה מִזְרָחָה׃}
}
{
	\eDtr1{
		and the plain, with the Jordan being the border, from the\\
		Kinneret to the sea of the plain, the Dead Sea, below the\\
		slopes of Pisgah eastward.\\
	}
}
{
}

\Verse{18}
{
	\hDtr1{וָאֲצַו אֶתְכֶם בָּעֵת הַהִוא לֵאמֹר יְהֹוָה אֱלֹהֵיכֶם נָתַן לָכֶם אֶת הָאָרֶץ הַזֹּאת לְרִשְׁתָּהּ חֲלוּצִים תַּעַבְרוּ לִפְנֵי אֲחֵיכֶם בְּנֵי יִשְׂרָאֵל כּל בְּנֵי חָיִל׃}
}
{
	\eDtr1{
		“And I commanded you at that time, saying, ‘YHWH, your\\
		God, has given you this land to take possession of it. You\\
		shall pass equipped, all the warriors, in front of your\\
		brothers, the children of Israel.\\
	}
}
{
}

\Verse{19}
{
	\hDtr1{רַק נְשֵׁיכֶם וְטַפְּכֶם וּמִקְנֵכֶם יָדַעְתִּי כִּי מִקְנֶה רַב לָכֶם יֵשְׁבוּ בְּעָרֵיכֶם אֲשֶׁר נָתַתִּי לָכֶם׃}
}
{
	\eDtr1{
		Only your wives and your infants and your livestock (I\\
		knew that you have a great amount of livestock) shall stay\\
		in your cities that I’ve given you\\
	}
}
{
}

\Verse{20}
{
	\hDtr1{עַד אֲשֶׁר יָנִיחַ יְהֹוָה לַאֲחֵיכֶם כָּכֶם וְיָרְשׁוּ גַם הֵם אֶת הָאָרֶץ אֲשֶׁר יְהֹוָה אֱלֹהֵיכֶם נֹתֵן לָהֶם בְּעֵבֶר הַיַּרְדֵּן וְשַׁבְתֶּם אִישׁ לִירֻשָּׁתוֹ אֲשֶׁר נָתַתִּי לָכֶם׃}
}
{
	\eDtr1{
		until YHWH will give your brothers rest like you, and\\
		they, too, will get possession of the land that YHWH, your\\
		God, is giving them across the Jordan. Then you shall go\\
		back, each to his possession that I’ve given you.’\\
	}
}
{
}

\Verse{21}
{
	\hDtr1{וְאֶת יְהוֹשׁוּעַ צִוֵּיתִי בָּעֵת הַהִוא לֵאמֹר עֵינֶיךָ הָרֹאֹת אֵת כּל אֲשֶׁר עָשָׂה יְהֹוָה אֱלֹהֵיכֶם לִשְׁנֵי הַמְּלָכִים הָאֵלֶּה כֵּן יַעֲשֶׂה יְהֹוָה לְכל הַמַּמְלָכוֹת אֲשֶׁר אַתָּה עֹבֵר שָׁמָּה׃}
}
{
	\eDtr1{
		And I commanded Joshua at that time, saying, ‘Your eyes,\\
		that have seen everything that YHWH, your God, has done to\\
		these two kings: so may YHWH do to all the kingdoms to\\
		which you’re crossing.\\
	}
}
{
}

\Verse{22}
{
	\hDtr1{לֹא תִּירָאוּם כִּי יְהֹוָה אֱלֹהֵיכֶם הוּא הַנִּלְחָם לָכֶם׃}
}
{
	\eDtr1{
		You shall not fear them, because YHWH, your God: He is the\\
		one fighting for you.’\\
	}
}
{
}

\Verse{23}
{
	\hDtr1{וָאֶתְחַנַּן אֶל יְהֹוָה בָּעֵת הַהִוא לֵאמֹר׃}
}
{
	\eDtr1{AND I IMPLORED “And I implored YHWH at that time, saying,}
}
{
%	\footnote{
%		I implored. Moses’ report of his anguished plea to God, begging to be
%		allowed to go to the promised land, is an exquisite reminder of another
%		report of anguish and begging: when Joseph’s brothers are accused of
%		spying in Egypt, they recall with shame and guilt that Joseph implored
%		them and they would not listen (Gen 42:21). Now Moses says that he
%		implored God, using the same word that Joseph’s brothers used—and these
%		are the only two occurrences of this word in the Torah—and that God
%		would not listen. And to cement the connection: Moses says that God told
%		him not to add anything more (“to go on speaking”) about this matter,
%		which in Hebrew is: Don’t tôsep—the root of the name Joseph! What is the
%		point of the parallel? When Joseph’s brothers do not listen to his plea,
%		he is the innocent one, and they are guilty. (They explicitly admit:
%		“But we’re guilty over our brother”—42:21.) But when Moses pleads, he is
%		not innocent. His fate has been sealed as a matter of justice. This is
%		extraordinary because when Moses pleads for the people in the matter of
%		the golden calf and the spies episode, he is successful: God forgives
%		them. There he is an innocent man pleading for mercy on behalf of the
%		guilty people. But now, as he pleads on his own behalf, he stands
%		guilty, and justice outweighs mercy. Another view: God is hardest on
%		those who are closest, like parents who may be tougher on their own
%		children than on the children of their friends. God shows no favoritism
%		for Moses. And the message is also that Moses’ serious public sin cannot
%		go unpunished. As God says in the matter of the death of Aaron’s sons:
%		“I shall be made holy through those who are close to me, and I shall be
%		honored in front of all the people” (Lev 10:3). Moses tells this to
%		Aaron on that occasion, and “Aaron was silent.” Now that message comes
%		back hauntingly to Moses, and God tells Moses that it is he who must be
%		silent. Thus it is here in Deuteronomy, the last book of the Torah, that
%		we see the Torah’s extraordinary embroidery. Words and themes wind, in
%		this case, from Genesis through Leviticus and Numbers into Deuteronomy,
%		like threads through a fabric.
%	}
}

\Verse{24}
{
	\hDtr1{אֲדֹנָי יֱהֹוִה אַתָּה הַחִלּוֹתָ לְהַרְאוֹת אֶת עַבְדְּךָ אֶת גּדְלְךָ וְאֶת יָדְךָ הַחֲזָקָה אֲשֶׁר מִי אֵל בַּשָּׁמַיִם וּבָאָרֶץ אֲשֶׁר יַעֲשֶׂה כְמַעֲשֶׂיךָ וְכִגְבוּרֹתֶךָ׃}
}
{
	\eDtr1{
		‘My Lord, YHWH, you’ve begun to show your servant your\\
		greatness and your strong hand, for who is a god in the\\
		skies and in the earth who can do anything like your acts\\
		and like your victories?!\\
	}
}
{
}

\Verse{25}
{
	\hDtr1{אֶעְבְּרָה נָּא וְאֶרְאֶה אֶת הָאָרֶץ הַטּוֹבָה אֲשֶׁר בְּעֵבֶר הַיַּרְדֵּן הָהָר הַטּוֹב הַזֶּה וְהַלְּבָנֹן׃}
}
{
	\eDtr1{
		Let me cross and see the good land that’s across the\\
		Jordan, this good hill country and the Lebanon.’\\
	}
}
{
%	\footnote{
%		Let me cross and see the good land that’s across the Jordan. Moses’ life
%		ends, as it began, at a river. On the continuing role that water plays
%		in his life, from the Nile to the Jordan, see the comment on Exod 17:6.
%		Recall also that the creation account in Genesis 1 begins with water,
%		and the account in Genesis 2 begins with rivers flowing out of Eden.
%		Moses’ life is thus a reflection in an individual human of the nature of
%		the world.
%	}
}

\Verse{26}
{
	\hDtr1{וַיִּתְעַבֵּר יְהֹוָה בִּי לְמַעַנְכֶם וְלֹא שָׁמַע אֵלָי וַיֹּאמֶר יְהֹוָה אֵלַי רַב לָךְ אַל תּוֹסֶף דַּבֵּר אֵלַי עוֹד בַּדָּבָר הַזֶּה׃}
}
{
	\eDtr1{
		“But YHWH was cross at me for your sakes, and He would not\\
		listen to me. And YHWH said to me, ‘You have much. Don’t\\
		go on speaking to me anymore of this thing.\\
	}
}
{
%	\footnote{
%		cross at me. This is not the usual word for anger. It occurs only here
%		in the Torah. It puns on the preceding verse: Moses says, “Let me cross
%		(Hebrew [image])” to the land that is “across ([image]) the Jordan.” And
%		then God is cross ([image]) and says “you shall not cross ([image]).”
%		The pun (in English and in Hebrew) conveys irony as well as the idea
%		that divine punishment in the Bible is often made to fit the crime.
%		Footnote 3:26. for your sakes. This has been taken to mean that Moses is
%		blaming the people for his fate. Alternatively, it can mean that God is
%		thinking of the people and is angry because Moses failed to sanctify God
%		in their eyes (Num 20:12). See also the comment on Deut 1:37. Footnote
%		3:26. You have much. These are the same words that Korah says to Moses
%		and Aaron. There is a great difference when the same words are spoken by
%		different persons—and with different motives. Korah uses them perhaps
%		out of envy, perhaps as a device to make his claim on priesthood sound
%		more justified. When God now uses them, it is more a reminder that Moses
%		has lived 120 years and has done great things. Even in being denied his
%		dream of coming to the land, he has much, perhaps more than any other
%		human who has ever lived.
%	}
}

\Verse{27}
{
	\hDtr1{עֲלֵה רֹאשׁ הַפִּסְגָּה וְשָׂא עֵינֶיךָ יָמָּה וְצָפֹנָה וְתֵימָנָה וּמִזְרָחָה וּרְאֵה בְעֵינֶיךָ כִּי לֹא תַעֲבֹר אֶת הַיַּרְדֵּן הַזֶּה׃}
}
{
	\eDtr1{
		Go up to the top of Pisgah and raise your eyes west and\\
		north and south and east and see it with your eyes,\\
		because you won’t cross this Jordan.\\
	}
}
{
}

\Verse{28}
{
	\hDtr1{וְצַו אֶת יְהוֹשֻׁעַ וְחַזְּקֵהוּ וְאַמְּצֵהוּ כִּי הוּא יַעֲבֹר לִפְנֵי הָעָם הַזֶּה וְהוּא יַנְחִיל אוֹתָם אֶת הָאָרֶץ אֲשֶׁר תִּרְאֶה׃}
}
{
	\eDtr1{
		And command Joshua and strengthen him and make him bold,\\
		because he will cross in front of this people, and he will\\
		get them the land that you’ll see as a legacy.’\\
	}
}
{
%	\footnote{
%		command Joshua. The story of the brothers’ sale of Joseph into slavery
%		again comes to mind, because that was the beginning of the process that
%		led the Israelites from their homeland to Egypt. And now, God’s response
%		to Moses ends with the charge to prepare Joshua to lead them, at last,
%		back into their homeland. And note that Joshua is from the tribe of
%		Ephraim; he is a descendant of Joseph! (And, like Joseph, he lives 110
%		years; Gen 50:26; Josh 24:29.)
%	}
}

\Verse{29}
{
	\hDtr1{וַנֵּשֶׁב בַּגָּיְא מוּל בֵּית פְּעוֹר׃}
}
{
	\eDtr1{“And we stayed in the valley opposite Bethpeor.}
}
{
}
